\zlabel{triptych:concise:themes:start}
\section{The Propers: Themes and Movement}

The Church begins Advent by asking the Son of God to come and save people who have reason to fear his judgment. At Mass they look forward both to the day when they will see their Judge and to the feast of his birth. What allows them to welcome him in either case is something he does for them first.

\subsection{The people who ask him to come}

The man whose words we sing at the entrance has already lifted his soul to God. His enemies can see whom he has trusted. Bellarmine asks us to consider what would happen if God were to fail him. His enemies could laugh at him for having placed all his hopes in someone who gave him no help. He takes heart from the others who have waited for God. None who wait for him will be put to shame. When he goes on to ask God to show him his paths, he has something more to learn than how to escape the men who despise him. In Bellarmine's account he wants to know how he should live to please God.\footnote{Ps 24(25):1--4. Robert Bellarmine, \work{Commentary on the Book of Psalms}, Ps 24, on vv. 1--4.}

At the Collect the priest speaks for everyone present. The danger he names arises from our sins. The prayer recalls Israel's appeal to its shepherd in Psalm 79(80). People who have turned away want their shepherd to turn them back and let them see his face. In the Collect the priest addresses the Son himself. The one who lives with the Father can come to protect those who have brought this danger on themselves and release them from it.\footnote{1962 Missal, First Sunday of Advent, no. 2. Ps 79:2--4 (80:1--3). Augustine, \work{Enarrationes in Psalmos}, English Psalm LXXX, §§2--3.}

Paul speaks to people who can remember first believing. Then they had good reason to change the way they lived. Why should they make the effort again now? Day is breaking. Paul tells them to wake from sleep and put on the armour of light. His last command in the lesson is to put on Christ. Quarrelling and envying someone else are among the things they must leave behind. In the Gradual the choir repeats the assurance about all who wait and asks God to make his paths known. Someone who heard Paul's list can ask to be taught how to behave toward the very person he envies.

The Alleluia asks for something the singers can see. They want God to show them his mercy and give them his salvation. Augustine asks what it is they have asked to receive. He answers that God's salvation is Christ. On his understanding, a person can look forward to the coming of his Judge for the same reason he looks forward to Christmas. He has asked to see the face of the one from whom he hopes for mercy.\footnote{Rom 13:11--14a. Ps 84:8 (85:7). Augustine, \work{Enarrationes}, English Psalm LXXXV, §6.}

\clearpage
\subsection{The one they expect and the one they receive}

In the Gospel there are people who would give much to be spared the sight before them. They hear the sea and are faint with fear at the thought of what will happen to the inhabited earth. Jesus tells his disciples to straighten up and lift their heads when these things begin. Gregory the Great takes the head to mean the mind that directs the body's members. Much that these men depended on may be taken away. They can still direct their attention to the person coming toward them. It is their redemption that Jesus says is near.\footnote{Lk 21:25--28. Gregory the Great, \work{Homiliae in Evangelia} I.1.1--3.}

Jesus has his disciples look at trees they have seen bear fruit. They know summer is coming before they have felt its heat. He tells them they can know the kingdom of God is near. In an earlier part of the same discourse he has warned them against following men who claim that the time has come. What the disciples have heard from Jesus about his own coming will remain true even when the heavens pass away. The return of the entrance antiphon at the Offertory lets someone who has heard this promise say before the gifts precisely whom he trusts.

The priest now prays to the Father through the Son. The holy things before him are to cleanse their recipients so that they can come with purer hearts to the one who gave them. In the dialogue before the Preface the people themselves answer the call to lift up their hearts. Once they have agreed that it is right to give thanks, the priest praises the Father for the divine majesty his Son and Spirit share. Those who heard of the Son appearing in majesty can already join the angels in the praise of God.\footnote{1962 Missal, First Sunday of Advent, nos. 7--8 and Preface of the Most Holy Trinity, no. 1082.}

At Communion the singers return to the psalm from which they asked to see God's mercy. They sing now of the goodness God will give and the fruit the earth will yield. The child born of Mary is that fruit for Bellarmine. Guéranger tells us that the choir sings this promise when the priest and people have received Communion. They still look ahead to the birth of the very Lord who has just given himself to them.\footnote{Ps 84:13 (85:12). Bellarmine, Ps 84, exposition 12. Prosper Guéranger, \work{The Liturgical Year}, Advent, First Sunday of Advent, p. 137.}

The Postcommunion is addressed by people who have received that gift and still have a favor to ask. Psalm 47(48) recalls mercy received in the temple. The priest changes the words of remembrance into a petition for mercy now. He asks it so that the worshippers can prepare to honor the coming festivals of their restoration. For a Christian expecting Christmas there is something more particular to hope for than the arrival of a day in the calendar. He can ask to keep the feast with the reverence due to the Son he has received.\footnote{Ps 47:10 (48:9). 1962 Missal, First Sunday of Advent, no. 10.}
\zlabel{triptych:concise:themes:end}
\clearpage
